\ifja
\chapter{時間軸の物理学：AssetとLiabilitiesの階層的分解}
\label{chap:time_horizon}

B/Sの左（資産）と右（負債）を無機質な勘定科目リストとして記憶してはならない。資産とは\textbf{「現金への変身カウントダウンタイマー」}であり、負債とは\textbf{「現金を奪う時限爆弾」}である。

\section{資産（Asset）の階層的分解式}
すべての資産は、現金化・回収までの「時間制限」によって次のように分解される。

\begin{equation}
    \text{Asset} = \text{Cash} + \text{NonCash}
\end{equation}
\begin{equation}
    \mathbf{\text{Asset} = \text{Cash} + \underbrace{\text{AR} + \text{Inv}}_{\ifja \text{1年流動（運転資本）}\else \text{1-Year Operational}\fi} + \underbrace{\text{Mid-Term}}_{\ifja \text{3年投資（IT/設備）}\else \text{3-Year Capital (IT/Equipment)}\fi} + \underbrace{\text{RealEstate} + \text{Intangibles}}_{\ifja \text{長期・固定（不動産/のれん）}\else \text{Illiquid Long-Term}\fi}}
\end{equation}
\begin{equation}
    \mathbf{\text{Liabilities} = \underbrace{\text{AP} + \text{Accrued}}_{\ifja \text{1年短期営業債務}\else \text{1-Year Operational}\fi} + \underbrace{\text{Short-Term Debt}}_{\ifja \text{1年内返済借入}\else \text{Due Within 1 Year}\fi} + \underbrace{\text{Mid-Term Debt}}_{\ifja \text{3年以内償還借入}\else \text{Due Within 3 Years}\fi} + \underbrace{\text{Long-Term Debt + Bonds}}_{\ifja \text{3年超長期負債}\else \text{Long-Term Liabilities}\fi}}
\end{equation}
\begin{equation}
\begin{aligned}
\ifja \text{左側（Asset: 現金化タイマー）} \else \text{Left (Asset: Liquidity Timer)} \fi &\quad \longleftrightarrow \quad \ifja \text{右側（Liabilities: 返済時限爆弾）} \else \text{Right (Liabilities: Time Bomb)} \fi \\
\text{Cash (0s)} &\quad \longleftrightarrow \quad \ifja \text{即時決済手形・未払金} \else \text{Immediate Obligations} \fi \\
\text{AR + Inv (1y)} &\quad \longleftrightarrow \quad \text{AP} + \ifja \text{短期借入 (1y)} \else \text{Short-Term Debt (1y)} \fi \\
\mathbf{\text{Mid-Term Assets (3y)}} &\quad \longleftrightarrow \quad \mathbf{\text{Mid-Term Debt (3y)}} \\
\text{Real Estate / Intangibles} &\quad \longleftrightarrow \quad \text{Long-Term Debt} + \text{Equity}
\end{aligned}
\end{equation}
\begin{align}
    \text{Net Debt} &\equiv (\text{Short-Term Debt} + \text{Mid-Term Debt}) - \text{Cash} \\
    \mathbf{\ifja \text{債務償還年数} \else \text{Debt Redemption Horizon} \fi} &= \mathbf{\frac{\text{Net Debt}}{\ifja \text{営業利益} + \text{減価償却費} \else \text{Operating Profit} + \text{Depreciation} \fi} \le 3.0 \quad [\ifja \text{年} \else \text{years} \fi]}
\end{align}


\begin{itemize}
    \item \textbf{Cash（即時：0秒）:} 究極の流動性。いかなる決済も拒絶されない基底状態。
    \item \textbf{短期1年流動（$\text{AR} + \text{Inv}$）:} 売掛金および在庫。日々の営業サイクルで1年以内に現金へ戻る運転資本。
    \item \textbf{中期3年投資（$\text{Mid-Term}$）：} PC、サーバー、小規模設備、内装、ソフトウェア。陳腐化が極めて速く、\textbf{3年以内の創出キャッシュで回収し尽くすべき}資産。
    \item \textbf{長期固定（$\text{RealEstate} + \text{Intangibles}$）：} 土地、建物、知財、M\&A時の買収プレミアム（のれん）。10年以上の超長期で拘束され、即座の換金性は絶望的である。
\end{itemize}

\section{負債（Liabilities）の階層的分解式}
負債側も全く同様に、返済期日（デッドライン）によって階層化される。

\section{時間軸マッチングの法則と「3年の壁」}
経営の安定性とは、左右のカウントダウンタイマーの速度が同期しているか否かで決まる。

もし左側が換金不能な不動産やのれんで埋まり、右側のタイマーが3年以内に集中していれば、たとえP/Lが黒字であろうと、時限爆弾が先に作動して会社は爆発（倒産）する。

\else
\chapter{Physics of Temporal Horizons: Hierarchical Decomposition of Assets and Liabilities}
\label{chap:time_horizon}

Assets are countdown timers to liquidity; liabilities are redemption time-bombs.

\section{Hierarchical Decomposition of Assets and Liabilities}
\begin{equation}
    \text{Asset} = \text{Cash} + \text{NonCash}
\end{equation}
\begin{equation}
    \mathbf{\text{Asset} = \text{Cash} + \underbrace{\text{AR} + \text{Inv}}_{\ifja \text{1年流動（運転資本）}\else \text{1-Year Operational}\fi} + \underbrace{\text{Mid-Term}}_{\ifja \text{3年投資（IT/設備）}\else \text{3-Year Capital (IT/Equipment)}\fi} + \underbrace{\text{RealEstate} + \text{Intangibles}}_{\ifja \text{長期・固定（不動産/のれん）}\else \text{Illiquid Long-Term}\fi}}
\end{equation}
\begin{equation}
    \mathbf{\text{Liabilities} = \underbrace{\text{AP} + \text{Accrued}}_{\ifja \text{1年短期営業債務}\else \text{1-Year Operational}\fi} + \underbrace{\text{Short-Term Debt}}_{\ifja \text{1年内返済借入}\else \text{Due Within 1 Year}\fi} + \underbrace{\text{Mid-Term Debt}}_{\ifja \text{3年以内償還借入}\else \text{Due Within 3 Years}\fi} + \underbrace{\text{Long-Term Debt + Bonds}}_{\ifja \text{3年超長期負債}\else \text{Long-Term Liabilities}\fi}}
\end{equation}
\begin{equation}
\begin{aligned}
\ifja \text{左側（Asset: 現金化タイマー）} \else \text{Left (Asset: Liquidity Timer)} \fi &\quad \longleftrightarrow \quad \ifja \text{右側（Liabilities: 返済時限爆弾）} \else \text{Right (Liabilities: Time Bomb)} \fi \\
\text{Cash (0s)} &\quad \longleftrightarrow \quad \ifja \text{即時決済手形・未払金} \else \text{Immediate Obligations} \fi \\
\text{AR + Inv (1y)} &\quad \longleftrightarrow \quad \text{AP} + \ifja \text{短期借入 (1y)} \else \text{Short-Term Debt (1y)} \fi \\
\mathbf{\text{Mid-Term Assets (3y)}} &\quad \longleftrightarrow \quad \mathbf{\text{Mid-Term Debt (3y)}} \\
\text{Real Estate / Intangibles} &\quad \longleftrightarrow \quad \text{Long-Term Debt} + \text{Equity}
\end{aligned}
\end{equation}
\begin{align}
    \text{Net Debt} &\equiv (\text{Short-Term Debt} + \text{Mid-Term Debt}) - \text{Cash} \\
    \mathbf{\ifja \text{債務償還年数} \else \text{Debt Redemption Horizon} \fi} &= \mathbf{\frac{\text{Net Debt}}{\ifja \text{営業利益} + \text{減価償却費} \else \text{Operating Profit} + \text{Depreciation} \fi} \le 3.0 \quad [\ifja \text{年} \else \text{years} \fi]}
\end{align}


\begin{itemize}
    \item \textbf{Cash (0s):} Absolute liquidity. Ground state of settlement.
    \item \textbf{Short-Term (1y):} Accounts receivable and inventory returning within normal operating cycles.
    \item \textbf{Mid-Term (3y):} IT infrastructure, equipment, tooling. Rapid obsolescence demands full recovery within 3 years.
    \item \textbf{Long-Term:} Real estate, patents, goodwill. Highly illiquid with multi-year exit frictions.
\end{itemize}

\section{The Three-Year Wall \& Debt Service Capability}
Dynamic stability demands synchronous matching of liquidity horizons against repayment deadlines. If obligations mature inside 3 years while assets are locked in illiquid holdings, structural default becomes mathematically inevitable regardless of P/L profitability.
\fi
